\section{Experimental Setup}

\paragraph{Reference decoder.}
The intended primary matrix pairs LLaDA-8B-Instruct and Dream-v0-Instruct-7B with
MATH-500 and GSM8K. Reference recipes must be pinned to upstream source and
validated independently before repairability is measured. An upstream benchmark
whose dataset, prompt, or evaluator differs from the paper setting is reported
with source-native smoke and recipe-fidelity status separate from the full
paper-task baseline; its reported score is not treated as directly comparable.
A different full official benchmark is not rerun solely as a prerequisite.

\paragraph{Reference trajectory bank and metric.}
Each backbone uses one sampled trajectory for each of 500 MATH-500 test items and
1,319 GSM8K test items. The main baseline is reference trajectory accuracy
(pass@1). The exact instrumented bank supplies failed and successful trajectory
pools; it is not regenerated for probing. Deep probing is restricted to a
deterministic failed-trajectory subset, ranked by a hash of design seed,
backbone, task, item identifier, and trajectory identifier. Selection never uses
repair outcomes.

\paragraph{Budget and priority.}
LLaDA targets 64 failed trajectories per task, with a minimum of 32. Dream
targets 32. The chosen count is frozen from measured
throughput, available GPU-hours, and remaining time before repair outcomes are
observed. LLaDA mechanism and temporal subsets each target 32 per task. LLaDA MATH deep
evidence is the critical path, followed by GSM8K replication. Dream validation
must not delay completion, sealing, or paper integration of LLaDA evidence.
These are design targets, not completed sample counts.

\paragraph{Mechanism and localization comparisons.}
The LLaDA mechanism comparison includes exact native replay, matched stochastic
continuation, canonical repair, random-position remasking with matched modified
position counts, CoRe-snapshot, and actually decoded fresh sampling at a matched
compute budget. CoRe-snapshot is a checkpoint adaptation, not a claim of
unmodified full-sequence CoRe reproduction. Localization reports confirmed
recovery, prospective trajectory-accuracy change, successful-trajectory harm,
coverage, the independently confirmed oracle gap, and NFE. The selector is not
replaced based on the observed reference result.

\paragraph{Measurement decoder.}
Existing V2 experiments use a separate measurement decoder and
\emph{measurement-decoder pass@8}. Their item-level pass@8 and all-eight-failed
item rate are distinct from reference pass@1 and the sampled-failed-trajectory
denominator. Existing LLaDA, Dream, BBH, and MBPP artifacts remain separate;
their eventual role as secondary evidence, decoder-regime robustness, or
appendix material depends on the sealed reference evidence.

\paragraph{Evidence boundary.}
Tables~\ref{tab:reference-existence}--\ref{tab:reference-localization} define the
reference reporting structure. Numerical cells require validated, sealed
reference bundles with artifact hashes and fixed execution provenance. An
em dash means that no sealed evidence has been integrated, not a zero estimate.
